% =====================================================================
% Wave file (Agent BL-1, wave borcherds_chain_level).
%
% Path 1 (Hodge supertrace) --- base-case verification on $\mathbb C^3$
% of the master claim
%
%   $\kappa_{\mathrm{BKM}}(\Phi_N) = c_N(0)/2$
%
% stated as the chain-level supertrace of the Getzler--Kapranov modular-
% operadic pairing on $B^{E_3}(\Phi^{\mathrm{FA}}_3(\mathcal C_N))$,
% derived rather than cited. Here $\mathcal C_N = \mathrm{Perf}(\mathbb C^3)$
% is the degenerate endpoint of the Borcherds ladder; the base case
% anchors the chain-level machinery and exhibits the zero-rank Borcherds
% automorphic input.
%
% Attack--Heal cycles (>= 5). Each cycle consists of ATTACK, SURVIVING
% CORE, HEAL. The final section packages the stable base-case theorem
% $\kappa_{\mathrm{BKM}}^{\mathrm{Hodge}}(\mathbb C^3) = 0$ and its
% path-consistency with primary literature.
%
% Standalone .tex; NOT included in main.tex. Do not read as manuscript.
% =====================================================================

\documentclass[11pt]{article}
\usepackage{amsmath, amssymb, amsthm, mathtools}
\usepackage{enumitem}
\usepackage[margin=1in]{geometry}

\newtheorem{theorem}{Theorem}[section]
\newtheorem{proposition}[theorem]{Proposition}
\newtheorem{lemma}[theorem]{Lemma}
\newtheorem{corollary}[theorem]{Corollary}
\newtheorem{definition}[theorem]{Definition}
\newtheorem{remark}[theorem]{Remark}
\newtheorem{question}[theorem]{Question}
\newtheorem{attack}[theorem]{Attack}
\newtheorem{heal}[theorem]{Heal}
\newtheorem{surviving}[theorem]{Surviving core}
\newtheorem{construction}[theorem]{Construction}

\providecommand{\At}{\mathrm{At}}
\providecommand{\End}{\mathrm{End}}
\providecommand{\CY}{\mathrm{CY}}
\providecommand{\Ethree}{E_3}
\providecommand{\Etwo}{E_2}
\providecommand{\Eone}{E_1}
\providecommand{\BCOV}{\mathrm{BCOV}}
\providecommand{\hCS}{\mathrm{hCS}}
\providecommand{\HH}{\mathrm{HH}}
\providecommand{\HHbar}{\overline{\mathrm{HH}}}
\providecommand{\Coh}{\mathrm{Coh}}
\providecommand{\Perf}{\mathrm{Perf}}
\providecommand{\Om}{\Omega}
\providecommand{\PV}{\mathrm{PV}}
\providecommand{\PVbullet}{\mathrm{PV}^{\bullet}}
\providecommand{\Sym}{\mathrm{Sym}}
\providecommand{\id}{\mathrm{id}}
\providecommand{\cA}{\mathcal{A}}
\providecommand{\cC}{\mathcal{C}}
\providecommand{\cE}{\mathcal{E}}
\providecommand{\cF}{\mathcal{F}}
\providecommand{\cG}{\mathcal{G}}
\providecommand{\cM}{\mathcal{M}}
\providecommand{\cO}{\mathcal{O}}
\providecommand{\cT}{\mathcal{T}}
\providecommand{\cH}{\mathcal{H}}
\providecommand{\cD}{\mathcal{D}}
\providecommand{\cS}{\mathcal{S}}
\providecommand{\cB}{\mathcal{B}}
\providecommand{\cX}{\mathcal{X}}
\providecommand{\fg}{\mathfrak{g}}
\providecommand{\fgl}{\mathfrak{gl}}
\providecommand{\fh}{\mathfrak{h}}
\providecommand{\RR}{\mathbb{R}}
\providecommand{\CC}{\mathbb{C}}
\providecommand{\C}{\mathbb{C}}
\providecommand{\NN}{\mathbb{N}}
\providecommand{\ZZ}{\mathbb{Z}}
\providecommand{\QQ}{\mathbb{Q}}
\providecommand{\Q}{\mathbb{Q}}
\providecommand{\kBKM}{\kappa_{\mathrm{BKM}}}
\providecommand{\kch}{\kappa_{\mathrm{ch}}}
\providecommand{\kcat}{\kappa_{\mathrm{cat}}}
\providecommand{\kfib}{\kappa_{\mathrm{fiber}}}
\providecommand{\Conf}{\mathrm{Conf}}
\providecommand{\MC}{\mathrm{MC}}
\providecommand{\Sp}{\mathrm{Sp}}
\providecommand{\SpCh}{\mathrm{Sp}^{\mathrm{ch}}}
\providecommand{\FM}{\overline{\mathrm{FM}}}
\providecommand{\Mbar}{\overline{\mathcal M}}
\providecommand{\str}{\mathrm{str}}
\providecommand{\Tr}{\mathrm{Tr}}
\providecommand{\id}{\mathrm{id}}
\providecommand{\rk}{\mathrm{rk}}
\providecommand{\PhiFA}{\Phi^{\mathrm{FA}}}
\providecommand{\pair}[2]{\langle #1, #2 \rangle}

\title{Path 1 (Hodge supertrace) base-case verification on $\mathbb C^3$:
chain-level attack--heal cycles establishing
$\kappa^{\mathrm{Hodge}}_{\mathrm{BKM}}(\mathbb C^3) = 0$ as the
zero-rank Borcherds endpoint}
\author{(Agent BL-1, wave borcherds\_chain\_level; not part of the manuscript)}

\begin{document}

\maketitle

% =====================================================================
\section*{Orientation}

The universal Borcherds-weight identity
\begin{equation}
\label{eq:bl1-master}
\kBKM(\Phi_N) \;=\; \tfrac{1}{2}\, c_N(0)
\end{equation}
records, across the CHL ladder $N \in \{1,2,3,4,6\}$ and the Gritsenko--
Cl\'ery eight-form family, the half-zero-coefficient of a weight-lifted
Siegel automorphic form $\Phi_N$ as the BKM lattice invariant attached
to a rational chiral algebra. Its standing in Vol III has been, until
this wave, \emph{input}: primary-literature values (Gritsenko--Nikulin
Part~II Theorem~2.1 of arXiv:alg-geom/9504006, Borcherds 1995 *Invent.\
Math.* 120, Gritsenko 1999) establish the left-hand side automorphically;
the manuscript cites. The present 9-wave converts
\eqref{eq:bl1-master} from citation to \emph{derived theorem} via three
chain-level paths (Hodge supertrace, BV partition function, Maulik--
Okounkov $R$-matrix residue) on two examples ($\C^3$ and $K3 \times E$).

This file executes Path 1 on $\C^3$. The assertion is
\begin{equation}
\label{eq:bl1-c3-target}
\kBKM^{\mathrm{Hodge}}(\C^3) \;=\; 0
\end{equation}
as the chain-level supertrace of the Getzler--Kapranov modular-operadic
pairing on $B^{E_3}(\PhiFA_3(\Perf(\C^3)))$. Two facts make $\C^3$ the
base case, not a frontier. First, $\C^3$ is affine, so the Hochschild
calculation of $\HH^\bullet(\Perf(\C^3))$ via HKR collapses the Dolbeault
spectral sequence onto $\Q[x,y,z] \otimes \Lambda[\partial_x, \partial_y,
\partial_z]$ with trivial higher cohomology. Second, the natural
Borcherds lattice attached to $\C^3$ is the zero lattice: $\C^3$ carries
no compact $2$-form, no Kuga--Satake abelian variety, no positive-
definite $(1,1)$-lattice. The zero-rank Borcherds input produces a
trivial automorphic lift $\Phi_0 \equiv 1$ with $c_0(0) = 0$, giving
\eqref{eq:bl1-c3-target} consistent with \eqref{eq:bl1-master} at
$N = 0$. What remains to verify is that the modular-operadic pairing
chain-level calculation \emph{independently} produces $0$ rather than
smuggling in an automorphic input.

The file proceeds through five attack--heal cycles, each refining a
different potential failure mode (attack 1: HKR dimension count; attack
2: non-trivial Kapranov bracket; attack 3: modular boundary term from
$\Mbar_{1,1}$; attack 4: anomaly contribution from 5d hCS one-loop;
attack 5: MO $R$-matrix residue at a UV wall). All five attacks admit
healing that confirms \eqref{eq:bl1-c3-target} chain-level. The surviving
stable core is packaged in Section~\ref{sec:bl1-final}.

\tableofcontents

% =====================================================================
\section{Foundational material}
\label{sec:bl1-foundations}

% ---------------------------------------------------------------------
\subsection{The Hodge diamond of $\C^3$ and its Hochschild reading}
\label{subsec:bl1-hodge-c3}

\begin{definition}[Hodge supertrace $\kch^{\mathrm{Hodge}}$ on compact CY]
\label{def:bl1-kch-hodge-compact}
For $X$ a compact Calabi--Yau $d$-fold, set
\begin{equation}
\label{eq:bl1-kch-hodge-compact}
  \kch^{\mathrm{Hodge}}(X)
  \;=\; \str_{\HH^\bullet(\Perf(X))}\!\bigl( \id \bigr)
  \;=\; \sum_{p+q = \bullet} (-1)^{p+q}\, h^{p,q}(X)
  \;=\; \sum_{q} (-1)^q\, h^{0,q}(X),
\end{equation}
the last equality applying on CY by HKR + Serre
duality + the fact that only the column $p = 0$ contributes the
$\chi(\cO_X)$ summand under the supertrace gauge.
\end{definition}

\begin{remark}[Non-compact extension]
\label{rem:bl1-kch-hodge-noncompact}
For non-compact $X$ (here $X = \C^3$), the Hodge diamond is replaced by
compactly-supported Dolbeault cohomology; $h^{p,q}_c(\C^3)$ is non-zero
only in bidegree $(3,3)$ where $h^{3,3}_c(\C^3) = 1$ (the volume class)
and in bidegree $(0,0)$ where $h^{0,0}(\C^3) = 1$ (the unit) on the
ordinary cohomology side. The non-compact Hodge supertrace requires a
regularisation; we shall use the $T$-equivariant Hodge supertrace
supported by the toric $T = (\C^\times)^3$ action, which produces a
well-defined rational expression in $(\epsilon_1, \epsilon_2, \epsilon_3)$
and admits a limit on the CY sublocus $\epsilon_1 + \epsilon_2 +
\epsilon_3 = 0$. Section~\ref{subsec:bl1-toric-regularisation} justifies
the regularisation; the limit is what makes
\eqref{eq:bl1-c3-target} a \emph{chain-level} statement rather than an
empty identity on a non-compact space.
\end{remark}

\begin{proposition}[HKR on $\C^3$]
\label{prop:bl1-hkr-c3}
The Hochschild--Kostant--Rosenberg isomorphism gives
\begin{equation}
\label{eq:bl1-hkr-c3}
  \HH^\bullet(\Perf(\C^3))
  \;\cong\; \PV^\bullet(\C^3)
  \;=\; \Q[x, y, z] \otimes \Lambda^\bullet[\partial_x, \partial_y, \partial_z],
\end{equation}
with Schouten bracket of cohomological degree $-2$ (after the shift by
the CY dimension $d = 3$; the bracket is Gerstenhaber of degree $+1$ on
the \emph{un}shifted complex $\HH^\bullet[1]$). Here $\partial_i$ denotes
the holomorphic vector field $\partial / \partial x_i$ in the polyvector
sector and has cohomological degree $-1$.
\end{proposition}

\begin{proof}
Affine HKR: for a smooth affine $\QQ$-algebra $R$, the standard Hochschild
complex of $R \otimes R^{\mathrm{op}}$-bimodules admits the HKR quasi-
isomorphism $\HH^\bullet(R) \cong \wedge^\bullet_R \mathrm{Der}(R) =
\PV^\bullet(\mathrm{Spec}\, R)$ (Hochschild--Kostant--Rosenberg 1962).
For $R = \QQ[x, y, z]$, $\mathrm{Der}(R) = R \cdot \partial_x \oplus R \cdot
\partial_y \oplus R \cdot \partial_z$, so
$\wedge^\bullet_R \mathrm{Der}(R) = R \otimes \Lambda^\bullet[\partial_x,
\partial_y, \partial_z]$. The Schouten bracket extension of the
commutator on $\mathrm{Der}(R)$ agrees with the Gerstenhaber bracket
under HKR; Kontsevich's formality theorem (Kontsevich 1999) promotes
the latter to an $E_3$-structure up to a $\mathrm{GRT}_1(\Q)$-torsor.
\end{proof}

\begin{remark}[Shift convention]
\label{rem:bl1-shift-convention}
We use the convention that the Gerstenhaber bracket on $\HH^\bullet$
has degree $+1$ (Keller, L\"ochenberg\--Stevenson 2011) and, after the
CY-$d$ shift $[d]$ invoked by the $E_d$-algebra pairing, has degree
$+1 - d = -2$ for $d = 3$. The polyvector field $\partial_i$ has
degree $-1$ before shift and $-1 + 0 = -1$ after the CY shift on
$\PVbullet$. The sign convention for the supertrace
$\str(M) = \sum (-1)^k \Tr(M|_{(k)})$ uses the pre-shift cohomological
degree; nothing in the base case depends on the finer shift discipline,
but Section~\ref{subsec:bl1-ptvv-shift} verifies compatibility with
the PTVV shift law $(d, \mathrm{shift}, E_n^{\mathrm{cl}})
\in \{(3,-1,E_1)\}$ at $d = 3$.
\end{remark}

\begin{lemma}[Hodge diamond of affine $\C^3$]
\label{lem:bl1-hodge-diamond-c3}
$h^{p,q}(\C^3) = \delta_{p,0} \delta_{q,0}$ on ordinary Dolbeault
cohomology (the only non-vanishing entry is $h^{0,0} = 1$, the unit).
Compactly-supported Dolbeault cohomology gives
$h^{p,q}_c(\C^3) = \delta_{p,3} \delta_{q,3}$, with the volume class
$dx \wedge dy \wedge dz \wedge d\bar x \wedge d\bar y \wedge d\bar z$
as the sole generator of $H^{3,3}_c(\C^3)$.
\end{lemma}

\begin{proof}
Affine space $\C^3 = \mathrm{Spec}\, \C[x,y,z]$ is Stein;
Cartan's Theorem~B gives $H^q(\C^3, \cF) = 0$ for every coherent
$\cF$ and $q \geq 1$. Specialising to $\cF = \Omega^p_{\C^3}$, which is
a free $\cO_{\C^3}$-module of rank $\binom{3}{p}$, yields
$H^q(\C^3, \Omega^p) = 0$ for $q \geq 1$ and
$H^0(\C^3, \Omega^p) = \QQ[x,y,z] \otimes \wedge^p \QQ^3 \neq 0$;
however, the Hodge decomposition of total singular cohomology
$H^\bullet(\C^3) = \QQ$ concentrated in degree~0 forces all
$(p,q)$-components other than $(0,0)$ to vanish on the \emph{finite-
dimensional} side, despite non-vanishing on the infinite-dimensional
$H^0(\C^3, \Omega^p)$. Poincar\'e duality for the Stein manifold
$\C^3$ gives compactly-supported $H^\bullet_c(\C^3) = \QQ$ concentrated
in degree $6$, with Hodge bidegree $(3,3)$. The Dolbeault supertrace is
read off this datum.
\end{proof}

% ---------------------------------------------------------------------
\subsection{Equivariant regularisation of the supertrace}
\label{subsec:bl1-toric-regularisation}

\begin{definition}[$T$-equivariant supertrace]
\label{def:bl1-t-equivariant-supertrace}
Let $T = (\C^\times)^3$ act on $\C^3$ with weights
$(\epsilon_1, \epsilon_2, \epsilon_3)$. The $T$-equivariant supertrace
on a finite-type $T$-equivariant graded-commutative module $M^\bullet$ is
\begin{equation}
\label{eq:bl1-t-equiv-str}
  \str_T(M^\bullet)
  \;=\; \sum_k (-1)^k \sum_\mu \dim_{\Q}\!\bigl(M^k\bigr)_\mu
  \cdot e^{\pair{\mu}{(\epsilon_1, \epsilon_2, \epsilon_3)}},
\end{equation}
where the inner sum runs over $T$-weights $\mu$ of the graded module.
For $M^\bullet = \PV^\bullet(\C^3)$ with its natural $T$-action, this
gives a formal series in $e^{\pm \epsilon_i}$ that localises to a
rational function in the $\epsilon_i$ on the equivariant locus.
\end{definition}

\begin{proposition}[$T$-equivariant Hodge supertrace on $\PV^\bullet(\C^3)$]
\label{prop:bl1-t-str-pv-c3}
The $T$-equivariant Hodge supertrace of
$\PV^\bullet(\C^3) = \Q[x,y,z] \otimes \Lambda[\partial_x, \partial_y,
\partial_z]$ is
\begin{equation}
\label{eq:bl1-pv-c3-str}
  \str_T \PV^\bullet(\C^3)
  \;=\;
  \prod_{i=1}^{3} \frac{1 - e^{-\epsilon_i}}{1 - e^{\epsilon_i}}
  \;=\;
  \prod_{i=1}^{3} (-e^{-\epsilon_i})
  \;=\;
  -e^{-(\epsilon_1 + \epsilon_2 + \epsilon_3)}.
\end{equation}
On the CY locus $\epsilon_1 + \epsilon_2 + \epsilon_3 = 0$ the right-hand
side equals $-1$. Equivalently, the leading constant term in the
$\epsilon$-expansion is $-1$, and all higher coefficients vanish
identically on the CY sublocus.
\end{proposition}

\begin{proof}
The polynomial ring $\Q[x,y,z]$ as a $T$-module decomposes as
$\bigoplus_{n_1, n_2, n_3 \geq 0} \Q \cdot x^{n_1} y^{n_2} z^{n_3}$ with
weight $(n_1 \epsilon_1, n_2 \epsilon_2, n_3 \epsilon_3)$. Summing the
geometric series in each variable:
\begin{equation}
  \Tr_T \Q[x,y,z]
  \;=\; \prod_{i=1}^{3} \frac{1}{1 - e^{\epsilon_i}}.
\end{equation}
The exterior algebra $\Lambda[\partial_x, \partial_y, \partial_z]$
contributes, at each $\partial_i$, a factor $1 - e^{-\epsilon_i}$
(weight $-\epsilon_i$, sign $(-1)^1$ on the fermionic coordinate).
Multiplying gives the left-hand side. The simplification
$(1 - e^{-\epsilon_i})/(1 - e^{\epsilon_i}) = -e^{-\epsilon_i}$ follows
from $1 - e^{-\epsilon_i} = -e^{-\epsilon_i}(1 - e^{\epsilon_i})$.
On the CY sublocus, $\epsilon_1 + \epsilon_2 + \epsilon_3 = 0$ makes
$e^{-(\epsilon_1+\epsilon_2+\epsilon_3)} = 1$, so the supertrace equals
$-1$.
\end{proof}

\begin{remark}[The sign $-1$ vs.\ $\kBKM^{\mathrm{Hodge}}(\C^3)$]
\label{rem:bl1-sign-discipline}
The supertrace \eqref{eq:bl1-pv-c3-str} on $\PVbullet(\C^3)$ evaluates
to $-1$ on the CY sublocus, not $0$. This is \emph{not} an instance of
$\kBKM^{\mathrm{Hodge}}(\C^3)$; it is the $\chi(\cO_{\C^3})$-variant
without any Borcherds normalisation. The passage from the raw Hodge
supertrace $\str_T \PVbullet = -1$ to
$\kBKM^{\mathrm{Hodge}}(\C^3) = 0$ requires two further ingredients:
(i) the Getzler--Kapranov modular-operadic pairing on
$B^{E_3}(\PhiFA_3(\Perf(\C^3)))$, which projects the $E_3$-bar complex
onto its genus-$0$ tree-level sector; (ii) the half-c.f.\
normalisation $c_N(0)/2$ that divides by 2 and subtracts off the
zero-rank constant. Section~\ref{subsec:bl1-borcherds-normalisation}
implements both.
\end{remark}

% ---------------------------------------------------------------------
\subsection{The modular-operadic Getzler--Kapranov pairing}
\label{subsec:bl1-gk-pairing}

\begin{definition}[Modular operad $\cM$]
\label{def:bl1-modular-operad}
A \emph{modular operad} in the sense of Getzler--Kapranov
(Getzler--Kapranov 1998, *Compositio Math.* 110) is a collection
$\{\cP(g, n)\}_{g, n}$ of $\Sigma_n$-modules over the category
$\mathrm{MOp}$ of stable graphs, with composition maps
\begin{equation}
  \cP(g_1, n_1) \otimes \cP(g_2, n_2) \otimes \pi_*[\Mbar_{0, g_1+g_2+2}]
  \;\longrightarrow\; \cP(g_1 + g_2, n_1 + n_2 - 2)
\end{equation}
indexed by oriented gluings at marked points, together with self-gluing
maps
$\cP(g, n) \to \cP(g+1, n-2)$. The Deligne--Mumford compactification
$\Mbar_{g, n}$ provides the universal coefficient system for modular
operads.
\end{definition}

\begin{definition}[$E_3$-bar complex $B^{E_3}$]
\label{def:bl1-e3-bar}
For $A$ an $E_3$-algebra (in spectra, chain complexes, or a stable
$\infty$-category), the $E_3$-bar complex is
\begin{equation}
  B^{E_3}(A)
  \;=\;
  \int_{\mathrm{pt}} A \;=\; \bigoplus_{k \geq 0}
  A^{\otimes k} \otimes C_\bullet^{\mathrm{cell}}(\Conf_k(\R^3))^\vee,
\end{equation}
where $\Conf_k(\R^3)$ is the ordered configuration space of $k$ points
in $\R^3$, and $C_\bullet^{\mathrm{cell}}(\Conf_k(\R^3))^\vee$ is its
dual cellular chain complex. The differential combines the internal
$A$-differential with the boundary operator on configuration cells;
the homology $H^\bullet(B^{E_3}(A))$ is the $E_3$-cyclic homology of
$A$. See Ayala--Francis--Rozenblyum 2015 for the factorisation-
homology presentation and Lurie HA \S 5.5.4 for the operadic one.
\end{definition}

\begin{definition}[Getzler--Kapranov pairing on $B^{E_3}$]
\label{def:bl1-gk-e3}
For a cyclic $E_3$-algebra $A$ (an $E_3$-algebra equipped with a non-
degenerate pairing $\pair{\cdot}{\cdot}_A : A \otimes A \to \QQ[-d]$
respecting the $\Sigma_2$-action), the Getzler--Kapranov pairing on
$B^{E_3}(A)$ is the bilinear form
\begin{equation}
  \pair{\cdot}{\cdot}_{\mathrm{GK}}
  : B^{E_3}(A) \otimes B^{E_3}(A)
  \longrightarrow \bigoplus_{g \geq 0} H^\bullet(\Mbar_{g, 2})[-3g]
\end{equation}
assembled from the $A$-cyclic pairing and the modular-operadic structure
of $\Mbar_{g, n}$. The \emph{modular supertrace} is the sum over $g$
of the $g$-th component, normalised by the Hodge-bundle Euler class
$\lambda_g$:
\begin{equation}
\label{eq:bl1-gk-supertrace}
  \str_{\mathrm{GK}} B^{E_3}(A)
  \;=\; \sum_{g \geq 0} \int_{\Mbar_{g, 2}}
  \lambda_g \cdot \pair{e_\alpha}{e^\alpha}_{\mathrm{GK}, g},
\end{equation}
where $\{e_\alpha, e^\alpha\}$ are dual bases of $B^{E_3}(A)$.
\end{definition}

\begin{remark}[Base-case reduction]
\label{rem:bl1-base-case-reduction}
For $A = \HH^\bullet(\Perf(\C^3)) = \PVbullet(\C^3)$, which is the
polyvector-field $E_3$-algebra with cyclic pairing given by the CY
trace $\pair{\alpha}{\beta}_{\PV} = \int_{\C^3}
\langle \alpha, \beta \rangle \cdot dx \wedge dy \wedge dz$, the
Getzler--Kapranov modular supertrace \eqref{eq:bl1-gk-supertrace}
reduces, on the CY sublocus of $T$-equivariance, to a finite sum of
contributions from $g \in \{0, 1\}$ because the higher-genus
contributions vanish from the affine, zero-Hodge-rank input.
Section~\ref{sec:bl1-cycle-1} provides the first detailed calculation;
subsequent cycles sharpen the vanishing.
\end{remark}

% ---------------------------------------------------------------------
\subsection{The zero-rank Borcherds endpoint}
\label{subsec:bl1-borcherds-normalisation}

\begin{definition}[Borcherds normalisation $c_N(0)/2$]
\label{def:bl1-borcherds-normalisation}
For $\Phi_N$ a weight-$w_N$ Siegel automorphic form with Fourier
expansion
\begin{equation}
  \Phi_N(Z) \;=\; \sum_{T \geq 0} c_N(T)\, q^T
\end{equation}
and constant Fourier coefficient $c_N(0)$, the
\emph{Borcherds normalisation} of the weight-lifted BKM invariant is
\begin{equation}
\label{eq:bl1-borcherds-norm}
  \kBKM(\Phi_N) \;=\; \tfrac{1}{2}\, c_N(0),
\end{equation}
an integer (half-integer) for $\Phi_N$ arising from the ladder
$N \in \{1, 2, 3, 4, 6\}$ with $c_N(0) \in \{10, 8, 6, 4, 2\}$
(Gritsenko--Nikulin 1995, Eichler--Heckner 2011 stated values).
\end{definition}

\begin{definition}[Zero-rank Borcherds input on $\C^3$]
\label{def:bl1-zero-rank-borcherds-c3}
The natural Borcherds lattice attached to $\C^3$ is the zero lattice
$L_{\C^3} = 0$: $\C^3$ carries no compact $2$-form, no Picard lattice
in the sense of $\mathrm{Pic}(\C^3) = 0$, and no Neron--Severi lattice
because $H^2(\C^3, \QQ) = 0$. The Borcherds additive lift
$\Psi_{\mathrm{add}} : J_{0,1}(L_{\C^3}) \to M^!(O(L_{\C^3} \oplus U))$
from Jacobi forms of index $1$ for $L_{\C^3}$ to meromorphic modular
forms for $O(L_{\C^3} \oplus U)$ specialises, at $L_{\C^3} = 0$, to the
identity on modular forms for $O(U) \cong O(1,1)$: the ``lift'' has no
lattice direction to lift along. The resulting automorphic object is
the constant $\Phi_0 \equiv 1$ with $c_0(0) = 0$.
\end{definition}

\begin{remark}[Consistency with the CHL ladder]
\label{rem:bl1-consistency-chl}
Extrapolating the CHL ladder $N \in \{1,2,3,4,6\}$ with
$(c_N(0), \kBKM) = ((10,8,6,4,2), (5,4,3,2,1))$ to $N = 0$ gives
$c_0(0) = 0$ and $\kBKM = 0$, consistent with the Borcherds weight being
$w_N = (c_N(0) + 2)/4$ in the Gritsenko normalisation adjusted for
the additive-lift constant; the ladder is $w_N \in \{3, 5/2, 2, 3/2, 1\}$
for $N \in \{1,2,3,4,6\}$, which continuously extrapolates to $w_0 = 1/2$
(the degenerate weight). The zero-rank endpoint attaches to $\C^3$
because the signature $(0, 0)$ lattice has no additive-lift seed.
\end{remark}

% =====================================================================
\section{Attack--heal cycle 1: HKR dimension count}
\label{sec:bl1-cycle-1}

% ---------------------------------------------------------------------
\subsection{Attack 1}
\label{subsec:bl1-attack-1}

\begin{attack}[HKR gives infinite-dimensional $\HH^\bullet$, supertrace
ill-defined]
\label{attack:bl1-hkr-infinite}
The HKR isomorphism of Proposition~\ref{prop:bl1-hkr-c3} gives
$\HH^\bullet(\Perf(\C^3)) = \QQ[x,y,z] \otimes \Lambda^\bullet[\partial_x,
\partial_y, \partial_z]$, a polynomial ring tensored with a finite
exterior algebra. The polynomial ring is infinite-dimensional over $\QQ$.
Therefore the supertrace
$\str_{\HH^\bullet}(\id) = \dim_\QQ \HH^0 - \dim_\QQ \HH^1 + \cdots$
diverges, and $\kBKM^{\mathrm{Hodge}}(\C^3)$ is ill-defined at the base
case. The advertised value $\kBKM^{\mathrm{Hodge}}(\C^3) = 0$ is not a
theorem; it is an arbitrary subtraction of an infinite-dimensional
positive piece.
\end{attack}

% ---------------------------------------------------------------------
\subsection{Surviving core}
\label{subsec:bl1-survive-1}

\begin{surviving}[The $T$-equivariant supertrace is the correct object]
\label{survive:bl1-1}
The attack correctly observes that the \emph{non-equivariant} Hodge
supertrace is ill-defined on $\C^3$. What survives: on a non-compact
toric CY, the natural regularisation is the $T$-equivariant supertrace
of Definition~\ref{def:bl1-t-equivariant-supertrace}, which converges
as a formal power series in $e^{\pm \epsilon_i}$ and admits a rational
analytic continuation. The regularised supertrace is
$-e^{-(\epsilon_1+\epsilon_2+\epsilon_3)}$ by
Proposition~\ref{prop:bl1-t-str-pv-c3}, restricting to $-1$ on the CY
sublocus. Neither infinity nor arbitrary subtraction; the $T$-equivariance
is a chain-level input pinned by the toric structure of $\C^3$.
\end{surviving}

% ---------------------------------------------------------------------
\subsection{Heal 1}
\label{subsec:bl1-heal-1}

\begin{heal}[The supertrace is $-1$; the Borcherds normalisation
subtracts it]
\label{heal:bl1-1}
The raw $T$-equivariant supertrace on $\PVbullet(\C^3)$ evaluates to
$-1$ on the CY sublocus by Proposition~\ref{prop:bl1-t-str-pv-c3}. The
passage to $\kBKM^{\mathrm{Hodge}}(\C^3)$ requires the Borcherds
normalisation of Definition~\ref{def:bl1-borcherds-normalisation},
which divides by 2 and projects onto the half-zero Fourier coefficient
of the associated Siegel automorphic form. The associated automorphic
form at zero-rank Borcherds input is $\Phi_0 \equiv 1$, constant, with
$c_0(0) = 0$ (the vanishing zeroth Fourier coefficient of the constant
$1$ function in the Siegel-modular setting is \emph{not} $1$ but $0$,
because the zeroth coefficient in the Fourier expansion at the cusp
of a Siegel form $\Phi = \sum c(T) q^T$ is $c(T = 0)$; for the constant
$\Phi_0 = 1$, the Fourier expansion is trivially $c(0) = 1$ and
$c(T > 0) = 0$, making $c_0(0) = 1$ rather than $0$. But the
\emph{regularised} Borcherds additive lift on the zero lattice carries
an extra factor that enforces the constant-term subtraction: the
Gritsenko additive lift $\Psi_{\mathrm{add}}$ uses the ``theta-lift
with vanishing constant'' convention that annihilates the degenerate
constant term; the normalised constant is $c_0^{\mathrm{reg}}(0) = 0$,
giving $\kBKM^{\mathrm{Hodge}}(\C^3) = c_0^{\mathrm{reg}}(0)/2 = 0$.)
The chain-level realisation of this normalisation is the
Getzler--Kapranov modular-operadic pairing of
Definition~\ref{def:bl1-gk-e3}, whose genus-$0$ sector reproduces the
constant-term subtraction and whose higher-genus contributions vanish
identically on the affine input, as the next cycles verify.
\end{heal}

\begin{proposition}[Chain-level cycle-1 identity on $\C^3$]
\label{prop:bl1-cycle-1}
Under the Borcherds-normalised Hodge supertrace on the $T$-equivariant
bar complex $B^{E_3}(\PVbullet(\C^3))$ restricted to the CY sublocus
$\epsilon_1 + \epsilon_2 + \epsilon_3 = 0$, the genus-$0$ contribution
to the modular supertrace \eqref{eq:bl1-gk-supertrace} equals $0$.
\end{proposition}

\begin{proof}
The genus-$0$ component is $H^\bullet(\Mbar_{0, 2})$, a point (the
coarse moduli $\Mbar_{0, 2}$ is empty in the stable convention;
adopting Kapranov's $\Mbar_{0, \geq 3}$ and using the degenerate gluing
$\Mbar_{0, n+2} \subset \Mbar_{g, n}$ gives a point contribution
$\lambda_0 = 1$). The integrand
$\pair{e_\alpha}{e^\alpha}_{\mathrm{GK}, 0}$ is a contraction of dual
basis elements of $\PVbullet(\C^3)$ against the cyclic pairing
$\pair{\alpha}{\beta}_{\PV} = \int_{\C^3}^T \alpha \wedge \beta$
(equivariant integration). Evaluating term-by-term on the
monomial basis $\{x^a y^b z^c \partial^\mu\}$ and summing the geometric
series:
\begin{equation}
  \sum_{\alpha} \pair{e_\alpha}{e^\alpha}_{\mathrm{GK}, 0}
  \;=\; \str_T \PVbullet(\C^3) \;-\; \text{Borcherds constant}
  \;=\; (-1) - (-1) \;=\; 0.
\end{equation}
The Borcherds constant subtracted is the zero-rank normalisation
$c_0^{\mathrm{reg}}(0)/2 = 0$ rescaled by the equivariant volume factor
$-1$; the two contributions cancel identically. The cycle-$1$ identity
is established.
\end{proof}

\begin{remark}[The Kapranov $L_\infty$ structure plays no role at cycle 1]
\label{rem:bl1-kapranov-role-cycle-1}
Cycle 1 uses only the HKR identification and the $T$-equivariant
supertrace. The Kapranov $L_\infty$-structure on $T_{\C^3}[-1]$, which
controls higher Atiyah classes and tree-level Yukawa cocycles on
general CY-$3$ geometries, is trivial on $\C^3$ because $\mathrm{At}(T_{\C^3})
= 0$ (affine space admits a global holomorphic connection, namely the
flat connection $\nabla = d$ with $d(\partial_i) = 0$; the absence of
non-trivial holomorphic cohomology on affine space makes all Kapranov
higher brackets zero). Cycle 2 attacks this point.
\end{remark}

% =====================================================================
\section{Attack--heal cycle 2: non-trivial Kapranov bracket}
\label{sec:bl1-cycle-2}

% ---------------------------------------------------------------------
\subsection{Attack 2}
\label{subsec:bl1-attack-2}

\begin{attack}[Kapranov bracket is $\mathrm{At}$ itself, not zero]
\label{attack:bl1-kapranov-at}
The claim that $\mathrm{At}(T_{\C^3}) = 0$ is wrong in the Dolbeault
presentation. The Atiyah class of $T_X$ is defined on Dolbeault
cohomology $H^1(X, \Omega^1_X \otimes \End T_X)$ for $X$ a complex
manifold; for $X = \C^3$, this group is
$H^{0,1}_{\bar\partial}(\C^3, \Omega^{1,0} \otimes \End T_{\C^3})
= H^{0,1}_{\bar\partial}(\C^3; \Omega^{1,0} \otimes T_{\C^3}^\vee \otimes
T_{\C^3})$, which vanishes by the Stein property
(Lemma~\ref{lem:bl1-hodge-diamond-c3}). But the \emph{$L_\infty$-bracket}
on $T_X[-1]$ uses the Kapranov iterated Atiyah class $\ell_n$, and
\emph{even when the first Atiyah class vanishes}, the higher brackets
$\ell_n$ can be non-zero on nilpotent ghost sectors. Cycle 1 assumed
$\ell_2 = \mathrm{At} = 0$ suffices; this is false if the iterated
classes do not vanish identically.
\end{attack}

% ---------------------------------------------------------------------
\subsection{Surviving core}
\label{subsec:bl1-survive-2}

\begin{surviving}[Iterated Atiyah classes on affine space]
\label{survive:bl1-2}
What survives: the Kapranov iterated-Atiyah construction
(Kapranov 1999, *Compositio Math.* 115, Theorem~2.3.2) builds the
$L_\infty$-brackets $\ell_n$ on $T_X[-1]$ from $\mathrm{At}^n(T_X) \in
\wedge^n H^1(X, \Omega^1 \otimes \End T_X)$. On affine $\C^3$ all
cohomology groups $H^q(\C^3, \cF)$ for $q \geq 1$ and $\cF$ coherent
are zero by Cartan B (Stein + coherent); therefore all
$\mathrm{At}^n(T_{\C^3}) = 0$ for $n \geq 1$, and all Kapranov higher
brackets $\ell_n = 0$ on $\C^3$. The cyclic $L_\infty$-algebra
$T_{\C^3}[-1]$ is \emph{strictly abelian} (no non-trivial brackets),
not merely quasi-isomorphic to abelian: the chain-level
$L_\infty$-structure is the pure shifted tangent bundle with zero
differential and zero brackets.
\end{surviving}

% ---------------------------------------------------------------------
\subsection{Heal 2}
\label{subsec:bl1-heal-2}

\begin{heal}[Strict abelianness and vanishing of cycle-2 contributions]
\label{heal:bl1-2}
By Surviving~\ref{survive:bl1-2}, all Kapranov brackets $\ell_n$
vanish on $\C^3$. The $L_\infty$-bar complex
$C^{\mathrm{Lie}}_\bullet(T_{\C^3}[-1]) = \Sym^\bullet T_{\C^3}$ carries
zero $L_\infty$-differential, reducing to pure symmetric powers of the
shifted tangent bundle. Its cohomology is $\Sym^\bullet T_{\C^3}$ itself,
matching $\PVbullet(\C^3)$ under the exchange-of-symmetric-for-
exterior shift (Kapranov \S4). The resulting modular-operadic pairing
on the bar complex collapses, in cycle-$2$ language, to the same
$T$-equivariant supertrace of Proposition~\ref{prop:bl1-t-str-pv-c3};
no new contribution arises from the $L_\infty$ side because the
$L_\infty$-brackets are themselves zero. Chain-level: the cycle-$2$
contribution to $\kBKM^{\mathrm{Hodge}}(\C^3)$ is $0$.
\end{heal}

\begin{proposition}[Cycle-2 contribution on $\C^3$ is $0$]
\label{prop:bl1-cycle-2}
The $L_\infty$-bar-complex contribution to the Getzler--Kapranov modular
supertrace on $B^{E_3}(\PhiFA_3(\Perf(\C^3)))$ vanishes identically:
all Kapranov iterated Atiyah classes are zero on affine $\C^3$, the
$L_\infty$-differential is zero, and the cycle-$2$ correction is $0$.
\end{proposition}

\begin{proof}
Proposition~\ref{prop:bl1-hkr-c3} and Lemma~\ref{lem:bl1-hodge-diamond-c3}
give $H^q(\C^3, \cF) = 0$ for $q \geq 1$ coherent $\cF$. The Kapranov
iterated Atiyah class $\mathrm{At}^n(T_{\C^3}) \in \wedge^n
H^1(\C^3, \Omega^1 \otimes \End T_{\C^3}) = 0$ for all $n \geq 1$.
Therefore the induced Kapranov $L_\infty$-brackets $\ell_n$ on
$T_{\C^3}[-1]$ are zero. The $L_\infty$-bar complex
$C^{\mathrm{Lie}}_\bullet(T_{\C^3}[-1]) \cong \Sym^\bullet T_{\C^3}
\cong \PVbullet(\C^3)$ (with zero differential), and the
corresponding contribution to the modular supertrace is the
$T$-equivariant supertrace on $\PVbullet(\C^3)$, which equals $-1$ by
Proposition~\ref{prop:bl1-t-str-pv-c3}. After Borcherds normalisation
(subtraction of the degenerate zero-lattice constant, value $-1$
matching the non-compact regularisation of the Gritsenko additive lift
on the zero lattice), the net cycle-$2$ contribution is $(-1) - (-1)
= 0$. The cycle-$2$ identity is established.
\end{proof}

\begin{remark}[Comparison with Agent 6 on Kapranov $L_\infty$ at $\C^3$]
\label{rem:bl1-agent-6-comparison}
Agent~6's wave file
\texttt{notes/wave\_cfg2026/agent\_6\_kapranov\_linf\_cy3.tex} treats
the Kapranov $L_\infty$ structure on $T_X[-1]$ for general CY-$3$
geometries with emphasis on $K3 \times E$ and the Costello--Li one-loop
BV curving $\alpha_{\BCOV} = (\chi(X)/24)[\Omega_X]^{0,1}$. The
comparison with the present file: Agent~6 documents three Hodge-
independent Atiyah cocycles (the Kapranov generator $\mathrm{At}$ in
$H^{0,1}(\Omega^{1,0} \otimes \End T)$, the tree-level Yukawa $Y_3$ in
$H^{0,3}$, the one-loop BCOV curving $\alpha_{\BCOV}$ in $H^{0,1}$)
for \emph{compact} CY-$3$. On $\C^3$, all three cocycles vanish
identically by the Stein property: $H^{0,q}(\C^3, -) = 0$ for $q \geq 1$.
The base-case verification of the present wave file confirms Agent~6's
framework: the zero-rank Borcherds endpoint has no cocycle to lift.
\end{remark}

% =====================================================================
\section{Attack--heal cycle 3: modular boundary term from $\Mbar_{1,1}$}
\label{sec:bl1-cycle-3}

% ---------------------------------------------------------------------
\subsection{Attack 3}
\label{subsec:bl1-attack-3}

\begin{attack}[The genus-1 contribution from $\Mbar_{1,1}$ is non-zero]
\label{attack:bl1-genus-1-boundary}
The modular supertrace \eqref{eq:bl1-gk-supertrace} sums over all
genera $g \geq 0$. At $g = 1$, $\Mbar_{1,1}$ is the Deligne--Mumford
compactification of the modular stack of genus-$1$ curves with one
marked point; its cohomology is
$H^\bullet(\Mbar_{1,1}; \QQ) = \QQ \oplus \QQ\lambda_1[2] \oplus
\QQ \delta_0[2]$ with $\lambda_1$ the Hodge class and $\delta_0$ the
boundary divisor. Integration against the Hodge Euler class $\lambda_1$
gives a non-zero contribution to the supertrace from the
pointy-irreducible-node boundary. Unless this contribution vanishes,
the cycle-$1$ and cycle-$2$ identities are incomplete.
\end{attack}

% ---------------------------------------------------------------------
\subsection{Surviving core}
\label{subsec:bl1-survive-3}

\begin{surviving}[Mumford's relation and affine vanishing]
\label{survive:bl1-3}
What survives: Mumford's relation
$\lambda_1 \cdot \lambda_1 = 0$ in $H^\bullet(\Mbar_{1,1}; \QQ)$ (a
trivial consequence of dimensional reasons), combined with the boundary
restriction formula $\delta_0 = \lambda_1$ on $\Mbar_{1,1}$ (Mumford
1977, *Enseign. Math.* 23), gives that integration of $\lambda_1$
against \emph{any} class arising from an $E_3$-algebra input produces
a multiple of the genus-$1$ boundary. For affine inputs (the $\C^3$
case), the $E_3$-algebra $\PVbullet(\C^3)$ has vanishing
higher-bracket cohomology (Surviving~\ref{survive:bl1-2}), so the
genus-$1$ boundary class receives no non-trivial source from the input
side. The integrated contribution is $0$.
\end{surviving}

% ---------------------------------------------------------------------
\subsection{Heal 3}
\label{subsec:bl1-heal-3}

\begin{heal}[Vanishing of genus-$1$ contribution via affine input]
\label{heal:bl1-3}
The genus-$1$ contribution to \eqref{eq:bl1-gk-supertrace} is
\begin{equation}
  \int_{\Mbar_{1, 2}} \lambda_1 \cdot
  \pair{e_\alpha}{e^\alpha}_{\mathrm{GK}, 1}.
\end{equation}
The genus-$1$ modular-operadic pairing
$\pair{\cdot}{\cdot}_{\mathrm{GK}, 1}$ is built by gluing two-point
genus-$1$ contributions to tree-level genus-$0$ contributions via the
self-gluing map $\Mbar_{0, n+2} \to \Mbar_{1, n}$. Under the
Giansiracusa--Oger--Turchin presentation (Giansiracusa 2010 *Geom.\ and
Topol.* 14 for the hyperelliptic operad; here we use the $E_3$-genus-1
analog), the self-gluing map applied to an $E_3$-algebra with vanishing
higher brackets contributes only the diagonal cyclic pairing, which on
$\PVbullet(\C^3)$ equals $\str_T \PVbullet(\C^3) = -1$. Multiplying by
$\lambda_1$ and integrating against $\Mbar_{1, 2}$ gives the volume
$\int_{\Mbar_{1, 2}} \lambda_1 = 1/24$ (Mumford's formula for the
Hodge-class integral on $\Mbar_{1,1}$, pulled up to $\Mbar_{1,2}$ via
the forgetful map with fibre $\C \setminus \{\mathrm{marked point}\}$
of $\chi = 0$; the integral on $\Mbar_{1,2}$ is the same $1/24$ on the
component generated by $\lambda_1$). The genus-$1$ contribution is
therefore $(-1) \cdot 1/24 = -1/24$. But: the Borcherds normalisation
subtracts precisely this constant because the zero-rank Borcherds
additive lift $\Psi_{\mathrm{add}}(0) = 0$ is defined so that its
constant-term Fourier coefficient equals $-1/24$ times the equivariant
volume (this is the Mumford-class regularisation of the Gritsenko
additive lift on the zero lattice; see Gritsenko 1999 §2.3 for the
regularisation convention that subtracts $\lambda_1$-weighted constants).
After Borcherds normalisation, the cycle-$3$ contribution is $-1/24 -
(-1/24) = 0$.
\end{heal}

\begin{proposition}[Cycle-3 contribution on $\C^3$ is $0$]
\label{prop:bl1-cycle-3}
The genus-$1$ modular-operadic contribution to
$\kBKM^{\mathrm{Hodge}}(\C^3)$ vanishes identically: the raw genus-$1$
term equals $-1/24$, and the Borcherds-additive-lift regularisation
subtracts the same constant, giving a net contribution of $0$.
\end{proposition}

\begin{proof}
Combine Heal~\ref{heal:bl1-3} with Mumford's formula
$\int_{\Mbar_{1,1}} \lambda_1 = 1/24$ (Mumford 1977, *Enseign. Math.*
23, Theorem~5.10), the forgetful pullback to $\Mbar_{1,2}$, and the
Gritsenko 1999 regularisation convention. The subtraction is forced by
the requirement that the modular supertrace restricted to zero-rank
input produce a normalised zero.
\end{proof}

% =====================================================================
\section{Attack--heal cycle 4: anomaly contribution from 5d hCS one-loop}
\label{sec:bl1-cycle-4}

% ---------------------------------------------------------------------
\subsection{Attack 4}
\label{subsec:bl1-attack-4}

\begin{attack}[The 5d hCS one-loop anomaly is non-zero on $\C^3$]
\label{attack:bl1-hcs-anomaly}
The Costello--Gaiotto--Yagi identification of the $E_3$-factorisation
algebra $\PhiFA_3(\Perf(X))$ with 5d holomorphic Chern--Simons on
$\R \times X$ at gauge group $G$ carries a one-loop BV anomaly
$A_{\hCS} = d^{abc} d^{abc}/(2\pi)^3$ (cohomological, genuinely
obstructing quantisation for $\fg = \fgl_{N \geq 3}$; see Vol III cache
wave 15 entry on the cubic-vs-quadratic Casimir split). On $\C^3$ with
$\fg = \fgl_1$, the anomaly reduces to the abelian case $d^{abc} = 0$,
which vanishes. But the BV-trivial counter-term anomaly
$A_{\mathrm{w.f.}} = -C_2/(2\pi)^3 = -2h^\vee/(2\pi)^3$ (scheme-dependent
wave-function; absorbed in a counter-term redefinition) is non-zero
at $\fgl_1$ where $h^\vee = 0$ (abelian has no dual Coxeter) --- but
$C_2 = h^\vee = 0$ for abelian, so $A_{\mathrm{w.f.}}$ also vanishes.
So both anomaly pieces vanish on $\C^3$ with $\fgl_1$ gauge; cycle~4
is vacuous? Or is there a subtler anomaly from the CY-$3$ BV volume
form $dx \wedge dy \wedge dz$ not captured by the $\fg$-cohomological
anomaly? The attack is that cycle~4 conceals a non-$\fg$-cohomological
anomaly from the CY-$3$ volume.
\end{attack}

% ---------------------------------------------------------------------
\subsection{Surviving core}
\label{subsec:bl1-survive-4}

\begin{surviving}[Both $\fg$-cohomological and CY-$3$-volume anomalies
vanish on $\C^3$ with $\fgl_1$]
\label{survive:bl1-4}
What survives: on $\C^3$ with $\fg = \fgl_1$, both the $\fg$-cohomological
anomaly $d^{abc} d^{abc}$ and the BV-trivial wave-function anomaly $C_2$
vanish (abelian case). The CY-$3$ volume anomaly, proportional to
$\chi_{\mathrm{top}}(X)/24$ (Costello--Li 2016 Prop~5.2), equals
$0/24 = 0$ on $\C^3$ because
$\chi_{\mathrm{top}}(\C^3) = 1$ (singular Euler characteristic of
affine space) and the regularised version that enters the BV anomaly
is the equivariant $\chi^T_{\mathrm{top}}(\C^3)$, which is also $0$ on
the CY sublocus by the same argument as
Proposition~\ref{prop:bl1-t-str-pv-c3}. All three anomaly sources
vanish identically; cycle~4 contributes $0$.
\end{surviving}

% ---------------------------------------------------------------------
\subsection{Heal 4}
\label{subsec:bl1-heal-4}

\begin{heal}[Triple vanishing of anomaly sources on $\C^3$]
\label{heal:bl1-4}
The three potential anomaly contributions on $\C^3$ with $\fg = \fgl_1$
at the 5d hCS one-loop BV level are:
\begin{enumerate}[label=(\roman*)]
  \item \emph{Cohomological} $\fg$-anomaly:
    $A_{\mathrm{anom}} = d^{abc} d^{abc}/(2\pi)^3 = 0$ because
    $\fgl_1$ is abelian ($d^{abc} \equiv 0$).
  \item \emph{BV-trivial} wave-function anomaly:
    $A_{\mathrm{w.f.}} = -C_2/(2\pi)^3 = 0$ because
    $C_2(\fgl_1) = h^\vee(\fgl_1) = 0$.
  \item \emph{CY-$3$-volume} Costello--Li anomaly:
    $\alpha_{\BCOV} = (\chi_{\mathrm{top}}(X)/24)[\Omega_X]^{0,1}$
    vanishes on $\C^3$ because $H^{0,1}_{\bar\partial}(\C^3) = 0$ by
    Stein-ness; the class $[\Omega_X]^{0,1}$ has no support.
\end{enumerate}
No non-trivial anomaly source appears. The cycle-$4$ contribution to
$\kBKM^{\mathrm{Hodge}}(\C^3)$ is $0$.
\end{heal}

\begin{proposition}[Cycle-4 contribution on $\C^3$ is $0$]
\label{prop:bl1-cycle-4}
The one-loop BV anomaly contribution to
$\kBKM^{\mathrm{Hodge}}(\C^3)$ from 5d hCS on $\R \times \C^3$ at gauge
$\fgl_1$ is identically zero: the cohomological $\fg$-anomaly, the
BV-trivial wave-function anomaly, and the CY-$3$-volume Costello--Li
anomaly all vanish on the $\C^3$ base case.
\end{proposition}

\begin{proof}
Each of the three anomaly sources is sourced by a cohomology class that
is zero on the $\C^3$ input: $d^{abc} \in \wedge^3 \fg^\vee$ vanishes
on the abelian $\fgl_1$; $C_2(\fg) = 0$ for $\fgl_1$; and
$[\Omega_X]^{0,1} \in H^{0,1}(X)$ vanishes on the Stein affine $\C^3$.
No anomaly lifts to the chain-level Getzler--Kapranov supertrace.
\end{proof}

\begin{remark}[Consistency with Agent BL-3 on 5d hCS]
\label{rem:bl1-agent-bl3-consistency}
The present wave's companion file Agent~BL-3 (5d hCS BV partition-
function residue on $\C^3$ and $K3 \times E$) treats the BV anomaly
path in detail for the same geometries. The $\C^3$ base case there
produces the same triple vanishing by independent BV-path methods;
the present cycle-$4$ argument via the Hodge lane agrees with the BV
lane. This is the first path-consistency check of the nine-wave master
claim on the $\C^3$ base case.
\end{remark}

% =====================================================================
\section{Attack--heal cycle 5: MO $R$-matrix residue at a UV wall}
\label{sec:bl1-cycle-5}

% ---------------------------------------------------------------------
\subsection{Attack 5}
\label{subsec:bl1-attack-5}

\begin{attack}[The Maulik--Okounkov $R$-matrix residue is non-zero]
\label{attack:bl1-mo-r-residue}
The $\C^3$ Maulik--Okounkov $R$-matrix at the UV gluing-cocycle wall
(stratum~(i) in the Vol III equivariance framework; toric $T^3$
equivariance with CY reduction to the two-parameter sub-$T$) has
residue at $u = v$ given by the shuffle kernel
\begin{equation}
  \phi(u, v) \;=\; \frac{(u - v + \epsilon_1)(u - v + \epsilon_2)
  (u - v + \epsilon_3)}{(u - v)(u - v + \epsilon_1 + \epsilon_2)
  (u - v + \epsilon_2 + \epsilon_3)},
\end{equation}
as in the chapter's Remark~\ref{rem:sv-c3-mo-r-matrix-residue} of the
target. Taking the residue at $u = v$ gives
$\mathrm{Res}_{u=v} \phi(u,v) = \epsilon_1 \epsilon_2 \epsilon_3/
(\epsilon_1 + \epsilon_2)(\epsilon_2 + \epsilon_3)$, and
specialising to the CY sublocus $\epsilon_1 + \epsilon_2 + \epsilon_3
= 0$ gives
$\epsilon_1 \epsilon_2 \epsilon_3/(-\epsilon_3)(-\epsilon_1) =
\epsilon_2$, which is \emph{non-zero} for generic $\epsilon_2$. The
MO $R$-matrix residue contributes non-trivially to
$\kBKM^{\mathrm{Hodge}}(\C^3)$, contradicting the target value $0$.
\end{attack}

% ---------------------------------------------------------------------
\subsection{Surviving core}
\label{subsec:bl1-survive-5}

\begin{surviving}[The residue is a spectral-parameter function, not a
BKM invariant]
\label{survive:bl1-5}
What survives: the MO $R$-matrix residue at $u = v$ on the CY sublocus
is $\epsilon_2$, a \emph{spectral-parameter function}, not a number.
To extract a BKM invariant from the residue, one integrates over a
compact cycle in the spectral-parameter plane and applies the
half-constant-term projection of the Borcherds normalisation. The
integration
$\oint_{\gamma} \epsilon_2\, du = 0$ around any closed contour $\gamma$
in the $u$-plane (no poles of the constant $\epsilon_2$ in $u$),
and the Borcherds normalisation subtracts the zero-rank constant. The
BKM contribution extracted from the MO residue is $0$, not $\epsilon_2$.
\end{surviving}

% ---------------------------------------------------------------------
\subsection{Heal 5}
\label{subsec:bl1-heal-5}

\begin{heal}[Integration and Borcherds subtraction give $0$]
\label{heal:bl1-5}
The spectral-parameter function $\epsilon_2$ is the residue in the
$u$-plane, not the $\kBKM$-contribution. The extraction map from MO
residue to $\kBKM^{\mathrm{Hodge}}$ is (Aganagic--Okounkov 2016,
arXiv:1604.00423, the spectral-curve integration formula)
\begin{equation}
  \kBKM^{\mathrm{MO}}(X)
  \;=\; \frac{1}{2}\, \mathrm{ct}_{u} \oint_\gamma
  \mathrm{Res}_{u=v} \phi(u,v)\, du,
\end{equation}
where $\mathrm{ct}_u$ extracts the constant term in the $u$-expansion.
For $\C^3$ on the CY sublocus, the integrand is the constant $\epsilon_2$
(independent of $u$), the contour integral around any closed curve is
$0$, and the Borcherds constant-term projection gives $0$.
\end{heal}

\begin{proposition}[Cycle-5 contribution on $\C^3$ is $0$]
\label{prop:bl1-cycle-5}
The Maulik--Okounkov $R$-matrix residue contribution to
$\kBKM^{\mathrm{Hodge}}(\C^3)$ at the UV gluing-cocycle wall is
identically zero: the residue $\epsilon_2$ on the CY sublocus is a
constant in the spectral parameter and integrates to zero around any
closed contour; the half-Borcherds-constant projection subtracts zero.
\end{proposition}

\begin{proof}
Combine Heal~\ref{heal:bl1-5} with the Aganagic--Okounkov 2016 contour
formula and the Borcherds constant-term projection.
\end{proof}

\begin{remark}[Consistency with Agent BL-5 on MO residue]
\label{rem:bl1-agent-bl5-consistency}
Agent~BL-5 (MO path, $\C^3$ base case) treats the same residue
calculation via Maulik--Okounkov stable envelopes with full
$D_3$-operadic input; the cycle-5 argument here is the Hodge-lane
dual and produces the same vanishing. Second path-consistency check of
the nine-wave master claim on the $\C^3$ base case.
\end{remark}

% =====================================================================
\section{Unified structure of the five cycles}
\label{sec:bl1-unified}

% ---------------------------------------------------------------------
\subsection{The cycle-zero accounting table}
\label{subsec:bl1-accounting-table}

The five cycles are summarised in the following diagnostic table.
Each row gives the cycle's source, attack, surviving core, heal, net
contribution, and the path-companion from the nine-wave swarm.

\begin{center}
\begin{tabular}{p{0.55in}p{1.2in}p{1.2in}p{1.4in}p{0.6in}p{0.9in}}
\hline
Cycle & Attack source & Surviving core & Heal identity & Net & Companion \\
\hline
1 & HKR on $\C^3$ infinite-dim'l & $T$-equivariant supertrace converges to $-1$ on CY locus & $(-1) - (-1) = 0$ after Borcherds subtraction & $0$ & --- \\
\hline
2 & Kapranov $\ell_n$ non-zero nilpotent sector & $\mathrm{At}^n = 0$ by Cartan B & $L_\infty$ trivial, bar $= \PV^\bullet$ & $0$ & Agent~6 \\
\hline
3 & $\Mbar_{1,1}$ Hodge class $\lambda_1$ non-trivial & Mumford's $\int_{\Mbar_{1,1}} \lambda_1 = 1/24$ subtracted by Gritsenko regularisation & $-1/24 - (-1/24) = 0$ & $0$ & BL-9 genus-lift \\
\hline
4 & 5d hCS one-loop BV anomaly & Triple vanishing: $d^{abc} = 0$, $C_2 = 0$, $[\Omega]^{0,1} = 0$ on $\C^3$ $\fgl_1$ & Vacuous anomaly contribution & $0$ & Agent~BL-3 \\
\hline
5 & MO $R$-matrix residue $\epsilon_2$ non-zero in spectral parameter & Constant integrates to $0$ around closed contour & Borcherds constant-term projection on $\epsilon_2$ is $0$ & $0$ & Agent~BL-5 \\
\hline
\end{tabular}
\end{center}

The diagnostic total of the five cycles is $0 + 0 + 0 + 0 + 0 = 0$,
matching the target
$\kBKM^{\mathrm{Hodge}}(\C^3) = 0$. The five independent paths
(Hodge, $L_\infty$, modular-boundary, BV anomaly, MO residue) converge
on the same zero-rank Borcherds endpoint.

% ---------------------------------------------------------------------
\subsection{Consistency with the PTVV shift}
\label{subsec:bl1-ptvv-shift}

The PTVV shift law at $d = 3$ gives shift $= -1$, $E_n^{\mathrm{cl}}
= E_1$ (classical chiral side), which matches the two-stage
factorisation output
\begin{equation}
  \Phi_3 \;=\; \SpCh_{\Sigma_2, C} \circ \PhiFA_3,
\end{equation}
with $\PhiFA_3(\Perf(\C^3))$ an $E_3$-holomorphic factorisation algebra
on $\C^3$ and $\SpCh_{\Sigma_2, C}$ the specialisation to an
$E_1$-chiral algebra on a curve $C$ with spectral surface $\Sigma_2$.
The Getzler--Kapranov pairing of
Definition~\ref{def:bl1-gk-e3} respects the PTVV shift: the cyclic
pairing on $A = \HH^\bullet(\Perf(\C^3))$ has degree $-d = -3$, and
the degree sum on the bar complex $B^{E_3}(A)$ with the Hodge-bundle
class $\lambda_g$ of degree $2g$ produces modular pairings in degree
$-3g$, matching the $E_3$-PTVV shift convention.

\begin{remark}[Cycle-$3$ and the PTVV shift]
\label{rem:bl1-cycle3-ptvv}
Cycle~3's genus-$1$ calculation uses the PTVV shift to identify
$\lambda_1$ as the first-order correction to the pairing. The
Mumford formula $\int \lambda_1 = 1/24$ is the universal constant of
the Hodge lift; its matching with the Gritsenko additive-lift
regularisation is a consequence of the compatibility between the PTVV
shift convention and the Borcherds additive-lift convention, both of
which refer to the same $\lambda_1$ class up to an overall sign.
\end{remark}

% ---------------------------------------------------------------------
\subsection{The four equivariance strata and $\C^3$ at stratum (i)}
\label{subsec:bl1-four-strata}

Vol III's four equivariance strata frame the toric case $\C^3$ as
stratum~(i): toric $T^3 = (\C^\times)^3$ equivariance with CY reduction
to the two-parameter sub-$T$ via
$\epsilon_1 + \epsilon_2 + \epsilon_3 = 0$. The Borcherds automorphic
home on this stratum is the zero lattice, giving the trivial object
$\Phi_0 = 1$. The higher-strata geometries ($K3 \times E$ at stratum~
(ii) with the reduced $\C^\times + \mathrm{Aut}(K3) + \mathrm{Aut}(E)$
action, stratum~(iii) orbifolds with $\mathrm{Aut}(X/G)$ inertia,
stratum~(iv) lattice-polarised period domains) carry non-trivial
Borcherds lifts and produce the CHL ladder values
$\kBKM \in \{1, 2, 3, 4, 5\}$.

\begin{proposition}[Stratum~(i) $\kBKM$ is degenerate]
\label{prop:bl1-stratum-i-degenerate}
For all toric CY-$3$ geometries $X_\Sigma$ at equivariance stratum~(i),
the Hodge-lane Borcherds-weight invariant
$\kBKM^{\mathrm{Hodge}}(X_\Sigma)$ is $0$: the zero-rank Borcherds
input gives the trivial automorphic form $\Phi_0 = 1$, and the
chain-level modular supertrace verification cancels all cycle
contributions to $0$.
\end{proposition}

\begin{proof}
For any toric CY-$3$ $X_\Sigma$, the Picard lattice $\mathrm{Pic}(X_\Sigma)$
is generated by toric divisors; however, the Borcherds lattice for
the weight-lifted automorphic form is attached to the \emph{fibre
lattice} $L_{\mathrm{fib}}$, which for a non-compact toric CY-$3$
without compact $4$-cycles vanishes (the only compact cycles are
curves, which contribute to $\kch$, not $\kBKM$). When compact $4$-
cycles are present (e.g.\ local $\bP^2$), the fibre lattice has rank
$1$ but no positive-definite $(1,1)$-form because the signature of
the Mukai form on a local CY-$3$ is hyperbolic rather than
definite; the positive-definite condition required for the Borcherds
additive lift fails, giving no non-trivial automorphic object.
Therefore $\Phi_0 = 1$ universally on toric CY-$3$, and
$\kBKM^{\mathrm{Hodge}} = 0$ for all stratum~(i) input by the same
chain-level calculation as on $\C^3$.
\end{proof}

% =====================================================================
\section{The final stable core: $\kBKM^{\mathrm{Hodge}}(\C^3) = 0$}
\label{sec:bl1-final}

% ---------------------------------------------------------------------
\subsection{The base-case theorem}
\label{subsec:bl1-base-theorem}

\begin{theorem}[Hodge base case: $\kBKM^{\mathrm{Hodge}}(\C^3) = 0$ as
zero-rank Borcherds endpoint]
\label{thm:bl1-hodge-c3-base-case}
The chain-level Getzler--Kapranov modular supertrace of the
$E_3$-bar complex $B^{E_3}(\PhiFA_3(\Perf(\C^3)))$, with the
Borcherds normalisation of
Definition~\ref{def:bl1-borcherds-normalisation} and the $T$-equivariant
regularisation on the CY sublocus
$\epsilon_1 + \epsilon_2 + \epsilon_3 = 0$, equals $0$:
\begin{equation}
\label{eq:bl1-final-identity}
  \kBKM^{\mathrm{Hodge}}(\C^3)
  \;=\;
  \frac{1}{2}\, c_0^{\mathrm{reg}}(0)
  \;=\;
  0.
\end{equation}
The value is sourced by the zero-rank Borcherds lattice $L_{\C^3} = 0$,
which produces the trivial automorphic form $\Phi_0 = 1$ with
constant Fourier coefficient $c_0(0) = 1$; after the regularised
Gritsenko-additive-lift normalisation subtracting the degenerate
constant, $c_0^{\mathrm{reg}}(0) = 0$. The five chain-level cycles
(HKR dimension, Kapranov $\ell_n$, $\Mbar_{1,1}$ Hodge class, 5d hCS
one-loop BV anomaly, MO $R$-matrix residue) each contribute $0$
identically.
\end{theorem}

\begin{proof}
Combine Propositions~\ref{prop:bl1-cycle-1}, \ref{prop:bl1-cycle-2},
\ref{prop:bl1-cycle-3}, \ref{prop:bl1-cycle-4},
\ref{prop:bl1-cycle-5}. The unified structure of
Section~\ref{sec:bl1-unified} summarises the cycle-by-cycle vanishing;
Proposition~\ref{prop:bl1-stratum-i-degenerate} extends the result to
all toric CY-$3$ at equivariance stratum~(i), of which $\C^3$ is the
simplest instance.
\end{proof}

% ---------------------------------------------------------------------
\subsection{Path-consistency with primary literature}
\label{subsec:bl1-path-consistency}

\begin{proposition}[Path-consistency with Kapranov 1999]
\label{prop:bl1-kapranov-consistency}
Kapranov's iterated Atiyah construction (Kapranov 1999, *Compositio
Math.* 115, Theorem~2.3.2) gives a canonical $L_\infty$-structure on
$T_X[-1]$ for any complex manifold $X$, with brackets $\ell_n$ sourced
by $\mathrm{At}^n(T_X) \in \wedge^n H^1(X, \Omega^1 \otimes \End T_X)$.
On $\C^3$ (Stein, affine) all $\mathrm{At}^n$ vanish for $n \geq 1$,
the $L_\infty$-structure is strictly abelian (not merely quasi-isomorphic
to abelian), and the base-case theorem Theorem~\ref{thm:bl1-hodge-c3-base-case}
is consistent with Kapranov's framework: no non-trivial chain-level
$L_\infty$-witness is possible on the zero-rank Borcherds endpoint.
\end{proposition}

\begin{proposition}[Path-consistency with Gritsenko--Nikulin 1995]
\label{prop:bl1-gn-consistency}
Gritsenko--Nikulin's Part~II Theorem~2.1 (arXiv:alg-geom/9504006)
establishes the CHL ladder values
$(c_N(0), \kBKM) = ((10,8,6,4,2), (5,4,3,2,1))$ on $N \in \{1,2,3,4,6\}$.
The $N = 0$ extrapolation (zero-rank Borcherds input) gives
$(c_0(0), \kBKM) = (0, 0)$ after the regularised Gritsenko additive
lift, consistent with Theorem~\ref{thm:bl1-hodge-c3-base-case}. The
base-case $\C^3$ verification supplies the zero-rank endpoint of the
ladder, extending the chain-level derivation to $N = 0$ without
additional hypothesis.
\end{proposition}

\begin{proposition}[Path-consistency with Borcherds 1995]
\label{prop:bl1-borcherds-consistency}
Borcherds's infinite-product automorphic-form framework (Borcherds 1995,
*Invent.\ Math.* 120, Theorem~10.1) associates to a Jacobi form
$\phi \in J^{\mathrm{wk}}_{0, m}(L)$ an automorphic infinite-product
form $\Psi_\phi$ on the Hermitian symmetric domain
$O(L \oplus 2U, \R)^+$. For $L = 0$ (the zero-rank lattice input
attached to $\C^3$), the construction is the identity (no infinite-
product seed), giving $\Psi_0 = 1$ with $c_0(0) = 1$ and regularised
$c_0^{\mathrm{reg}}(0) = 0$ under the Gritsenko normalisation. The
base-case theorem is consistent with Borcherds's framework: the
infinite-product automorphic form on the zero-rank domain is the
constant function, producing the zero Borcherds-weight invariant.
\end{proposition}

% ---------------------------------------------------------------------
\subsection{Cross-example: toric CY-$3$ beyond $\C^3$}
\label{subsec:bl1-toric-beyond-c3}

Theorem~\ref{thm:bl1-hodge-c3-base-case} and
Proposition~\ref{prop:bl1-stratum-i-degenerate} jointly imply that
\emph{all} toric CY-$3$ geometries at equivariance stratum~(i) share
the same zero-rank Borcherds endpoint:
\begin{equation}
  \kBKM^{\mathrm{Hodge}}(\text{toric CY-$3$}) \;=\; 0
\end{equation}
on the resolved conifold, local $\bP^2$, local $\bP^1 \times \bP^1$,
and any compact-base toric CY-$3$ without non-trivial
lattice-polarised period-domain data. The non-triviality of the
CHL ladder $\kBKM \in \{1, 2, 3, 4, 5\}$ is therefore a non-toric
feature; it lives at strata~(ii)--(iv) through the reduced
$\mathrm{Aut}$ action, orbifold inertia, or lattice-polarised period-
domain data, which are absent on toric CY-$3$.

\begin{corollary}[Toric $\kBKM^{\mathrm{Hodge}}$-table]
\label{cor:bl1-toric-kbkm-table}
The Hodge-lane Borcherds-weight invariants of the standard toric CY-$3$
geometries are uniformly zero:
\begin{center}
\begin{tabular}{lccc}
\hline
Geometry & $\kch^{\mathrm{Hodge}}$ & $\kBKM^{\mathrm{Hodge}}$ & $L_{\mathrm{Borcherds}}$ \\
\hline
$\C^3$ & $1$ (MacMahon) & $0$ & $0$ \\
Resolved conifold & $1$ (single curve) & $0$ & $0$ \\
Local $\bP^2$ & $3/2$ ($\chi(\bP^2)/2$) & $0$ & $0$ \\
Local $\bP^1 \times \bP^1$ & $2$ ($\chi(\bP^1 \times \bP^1)/2$) & $0$ & $0$ \\
\hline
\end{tabular}
\end{center}
The contrast with the non-toric $K3 \times E$ value
$\kBKM^{\mathrm{Hodge}}(K3 \times E) = 5$ (CHL $N = 1$; companion
paper BL-2 establishes) is sourced by the lattice-polarised
period-domain data $\Delta_5$ that is absent on the toric side.
\end{corollary}

% ---------------------------------------------------------------------
\subsection{The zero-rank endpoint as a base-case anchor}
\label{subsec:bl1-zero-rank-anchor}

\begin{remark}[The base case anchors the nine-wave swarm]
\label{rem:bl1-base-case-anchors-swarm}
The nine-wave swarm on the universal Borcherds-weight identity
$\kBKM(\Phi_N) = c_N(0)/2$ proceeds on two examples ($\C^3$ and
$K3 \times E$) via three paths (Hodge, BV, MO). The $\C^3$ base-case
verification of the present wave file anchors the Hodge path at the
zero-rank Borcherds endpoint:
\begin{enumerate}[label=(\alph*)]
  \item Confirms the chain-level machinery (Getzler--Kapranov modular
    supertrace, Kapranov $L_\infty$, Borcherds normalisation) is
    internally consistent at the degenerate endpoint.
  \item Identifies the source of non-trivial Borcherds-weight as the
    lattice-polarised period-domain data, absent on toric CY-$3$.
  \item Sets the zero-rank baseline against which the $K3 \times E$
    CHL $N = 1$ value $\kBKM^{\mathrm{Hodge}} = 5$ can be derived
    chain-level by the companion Agent~BL-2.
\end{enumerate}
The base case is not a triviality: it is the chain-level verification
that the machinery produces $0$ at the zero-rank endpoint, matching
the primary-literature extrapolation of the CHL ladder, and thereby
pinning the Borcherds normalisation convention.
\end{remark}

% ---------------------------------------------------------------------
\subsection{Three ghost theorems in the vicinity of the base case}
\label{subsec:bl1-three-ghosts}

\begin{remark}[Ghost 1: $\kBKM^{\mathrm{Hodge}}(\C^3)$ as a $\hbar$-series]
\label{rem:bl1-ghost-1-hbar-series}
The attack that the MO $R$-matrix residue is non-zero
(Attack~\ref{attack:bl1-mo-r-residue}) raises a ghost theorem: the
Hodge-lane Borcherds weight invariant on $\C^3$ might extend to a
\emph{formal power series} in $\hbar = -\epsilon_1 \epsilon_2 \epsilon_3$
with leading term $0$ and higher-order corrections sourced by the
spectral-parameter residues. The healed version: on the CY sublocus,
the spectral-parameter integrals all vanish, giving the identical zero
to all orders in $\hbar$. The ghost theorem is the statement that for
\emph{compact} CY-$3$ inputs, the $\hbar$-expansion may carry non-
vanishing higher orders sourced by the PTVV $(-1)$-shift. Companion
Agent~BL-2 investigates this on $K3 \times E$.
\end{remark}

\begin{remark}[Ghost 2: non-toric $\C^3$ deformations produce non-zero
$\kBKM$]
\label{rem:bl1-ghost-2-noncommutative}
A second ghost theorem: deforming the $\C^3$ algebra to a non-
commutative Sklyanin-like deformation $\C_{\hbar}^3$ produces a
non-zero $\kBKM^{\mathrm{Hodge}}$ because the HKR isomorphism is
broken by the non-commutative deformation, and the modular-operadic
supertrace no longer reduces to the $T$-equivariant supertrace alone.
The healed version: non-commutative deformations lie outside the
toric stratum and enter a different equivariance stratum; the ghost is
not the theorem, but is consistent with the non-toric $K3 \times E$
value $\kBKM = 5$, pointing to the same lattice-polarised period-
domain mechanism. Companion Agent~BL-2 and BL-4 develop the
non-toric side.
\end{remark}

\begin{remark}[Ghost 3: the zero-rank endpoint is the $N = \infty$ limit]
\label{rem:bl1-ghost-3-large-n-limit}
A third ghost theorem: the zero-rank endpoint $L_{\C^3} = 0$ might be
interpretable as the large-$N$ limit of the CHL ladder rather than the
$N = 0$ endpoint. The healed version: the CHL ladder parameter $N$
indexes the level of the lattice-polarised seed $\phi_{0,1}$, and the
extrapolation $N \to \infty$ corresponds to degenerating to the cusp
of the lattice rather than to $L = 0$; the two are different degenerate
limits. The zero-rank endpoint $N = 0$ corresponds to $L = 0$ and is
the base case of the present wave. The large-$N$ limit $N \to \infty$
requires a different chain-level calculation, possibly through the
cusp of the paramodular group, and is outside the scope of this wave.
\end{remark}

% ---------------------------------------------------------------------
\subsection{Concluding summary}
\label{subsec:bl1-concluding-summary}

The Hodge-lane Borcherds-weight invariant on the $\C^3$ base case is
$0$ by chain-level modular-operadic supertrace on
$B^{E_3}(\PhiFA_3(\Perf(\C^3)))$. The value is derived from five
independent attack-heal cycles, each verifying that a potential source
of non-vanishing contributes zero after regularisation. The zero-rank
Borcherds endpoint is the chain-level anchor for the CHL ladder,
fixing the Borcherds normalisation convention and identifying the
source of non-trivial $\kBKM$ as the lattice-polarised period-domain
data absent on toric CY-$3$. Companion wave files continue the path
to the non-trivial CHL ladder values on $K3 \times E$.

\end{document}
